\chapter{The Algebra of Records}
\label{chap:algebra-of-records}

% Closed question: How do multiple records combine without collapsing into one
% undifferentiated record?

\begin{chapteressay}

A matrix $A$ has dimensions: a row count $m$ and a column count $n$.
Two matrices $A$ and $B$ can be multiplied to form $AB$ only when the
column count of $A$ matches the row count of $B$. If $A$ is $3 \times
4$ and $B$ is $4 \times 2$, the product $AB$ is well-defined and is
$3 \times 2$. If $A$ is $3 \times 4$ and $B$ is $5 \times 2$, the
product is undefined; the dimensions are incompatible.

This is a small algebraic fact about matrices. It is also the
chapter's organizing observation. Two records can be combined only
when their structural types align. The types specify the
compatibility conditions; the compatibility conditions determine the
admissible operations.

The chapter is about the algebra of records: how multiple records
combine into structured composites without losing their individual
identities. The chapter develops five mathematical structures.

The first is compatibility. Two records can be combined into a
product only when their types align. Matrix multiplication requires
matching column-row dimensions. Dimensional analysis requires
matching units. Function composition requires matching domain and
codomain. In each case, compatibility is a structural condition;
the operation is admissible only when the condition is met.

The second is product structure. The Cartesian product of two sets
$A$ and $B$ is the set of pairs $(a, b)$ with $a \in A$ and $b \in
B$. The product preserves the source identity of each component: the
projections $\pi_A$ and $\pi_B$ recover the components. Multi-
component records combine via products without losing the components.

The third is source/encode/execute. A mathematical record progresses
through stages: source (abstract specification), encoded
(materialized intermediate), executed (final value). The stages are
distinct; the transformations between stages are admissible
operations. The TeX-to-PDF-to-screen example makes this concrete:
the TeX file is the source; the PDF is the encoded form; the
rendered screen is the executed output.

The fourth is unresolved states. A record may carry components that
are admissible but not yet resolved. Formal power series carry an
unresolved variable $x$; they admit operations (addition,
multiplication, differentiation) without resolving the variable.
The chapter's lesson is that unresolved states are admissible; they
are not errors or defects.

The fifth is staging. Records have stages corresponding to their
progressive realization. The lambda calculus expression $(\lambda x. x +
1) \, 3$ stages through beta reduction ($3 + 1$) and arithmetic
evaluation ($4$). Each stage is a different mathematical object;
the algebra tracks the transformations explicitly.

The chapter's ground example is matrix algebra. The chapter's
second example is the stage production before opening night. A
theatrical production has five distinct records: the script, the
blocking, the lighting cue sheet, the sound cue sheet, and the
rehearsal log. Each record is mathematically distinct; the records
are related by compatibility constraints; the production is the
product of the records; the production stages through the
rehearsal process from read-through to opening night.

The chapter's closed question is:

\begin{center}
\emph{How do multiple records combine without collapsing into one
undifferentiated record?}
\end{center}

The chapter's answer is: the algebra of records. Compatibility
determines admissible combinations; product structure preserves
source identity; source/encode/execute tracks the stages of
realization; unresolved states are admissible; staging tracks the
transformations.

Subsequent chapters build on this algebra. Chapter 5 introduces
distance from order: the gap between recorded marks becomes an
admissible distance through the count of intermediate
distinguishable values. Chapter 6 introduces motion: the
selection of one continuous trajectory from many admissible ones.
Chapter 7 introduces transport: how information moves between
records while preserving the algebra.

The chapter is foundational rather than load-bearing. Its
structures are simple; they are the algebraic prerequisites for
the more elaborate constructions in later chapters. Without
compatibility, products, staging, and unresolved states, the
later chapters would lack the algebraic vocabulary for their
claims. The chapter establishes the vocabulary.

\end{chapteressay}

\section{Compatible Records}
\label{se:compatible-records}

A matrix $A \in \mathbb{R}^{m \times n}$ is a rectangular array of real
numbers with $m$ rows and $n$ columns. Two matrices $A$ and $B$ can be
multiplied to form $AB$ only when the number of columns of $A$ equals
the number of rows of $B$. If $A$ is $3 \times 4$ and $B$ is $4 \times
2$, then $AB$ is well-defined and is $3 \times 2$. If $A$ is $3
\times 4$ and $B$ is $5 \times 2$, then $AB$ is undefined: the
dimensions are incompatible.

This compatibility condition is not arbitrary. It is a structural
consequence of what matrix multiplication is. The product $AB$ is
defined entrywise: $(AB)_{ik} = \sum_j A_{ij} B_{jk}$. The index $j$
runs over the columns of $A$, which must equal the rows of $B$ for
the sum to make sense. If the dimensions do not match, the sum has
no meaning; the product is not defined.

The chapter's first claim is that compatibility is the algebraic
prerequisite for combination. Two records can be combined into a
larger structure only when they satisfy a compatibility condition.
The condition is part of the records' types, not an external
restriction. A record's type specifies what it can be combined with;
records with incompatible types cannot be combined.

The matrix example is the canonical instance. The matrix type
$\mathbb{R}^{m \times n}$ carries two parameters: the row count and
the column count. Two matrices are compatible for multiplication
when one's column count equals the other's row count. The product
inherits the row count of the first and the column count of the
second.

The compatibility extends to other operations. Two matrices of the
same shape can be added: $A + B$ requires both to be $m \times n$
with the same $m$ and $n$. A scalar can be multiplied by any
matrix: $cA$ requires $A$ to be a matrix and $c$ to be a scalar
(specifically, an element of the underlying field). Each operation
has its own compatibility condition; the conditions are encoded in
the operation's type signature.

In Lean, compatibility is enforced at the type level. The
typeclasses \texttt{Add} and \texttt{Mul} require their operands
to be of specific types. Matrix multiplication is a function
\texttt{Matrix.mul : Matrix m n R $\to$ Matrix n p R $\to$ Matrix
m p R}: the input types specify the compatibility ($n$ appears in
both), and the output type is determined.

The compatibility condition is not just about types. It can also
be a runtime check. Two physical measurements with different
units are typically incompatible: adding meters to seconds is not
admissible, even though both are real numbers. The dimensional
analysis of physics is the runtime version of compatibility: each
quantity carries its unit, and combinations are admissible only
when the units agree.

\begin{phenom}{Dimensional Consistency}
\label{ph:dimensional-consistency}

\PhStatement Quantities with units can be added or compared only when
their units agree. The product of two quantities with units $[u_1]$
and $[u_2]$ has unit $[u_1 \cdot u_2]$; the quotient has unit $[u_1 /
u_2]$.

\PhOrigin Dimensional analysis as a discipline emerged in the
nineteenth century, formalized by Fourier in the theory of heat. It
is the basis of all modern scientific units (SI, CGS, imperial).

\PhObservation The expression $5 \text{ m} + 3 \text{ s}$ is undefined.
The expression $5 \text{ m} \cdot 3 \text{ s} = 15 \text{ m s}$ is
defined and has units of meter-seconds. The expression $5 \text{ m}
/ 3 \text{ s} \approx 1.67 \text{ m/s}$ is defined and has units of
meters per second (velocity).

\PhConstraint Two scientific records can be combined only when
their dimensional structures admit the combination. Records with
incompatible dimensions cannot be added or subtracted; their
multiplication or division produces a new dimension.

\PhConsequence Compatibility is not an aesthetic preference; it is
a structural requirement of the algebra of records. Software
libraries that enforce dimensional consistency (e.g.,
\texttt{Units.jl}, \texttt{boost::units}) implement the algebra at
the type level, refusing to compile dimensionally inconsistent
operations.

\end{phenom}

The Dimensional Consistency phenomenon shows that compatibility
extends from pure mathematics into the structure of scientific
measurement. Records are not just numbers; they are numbers with
types. The types include dimensional information. Combinations
admissible at the level of bare numbers may not be admissible at
the level of typed records.

The chapter's claim is that algebra is the structure of admissible
combinations. The algebraic operations (sum, product, composition)
have type signatures that specify their compatibility conditions. A
record is a typed entity; combinations of records are admissible
only when the type signatures align. The discipline is that
admissible operations are derived from the type signatures, not
from the underlying numerical content.

\section{Product Structure}
\label{se:product-structure}

The Cartesian product of two sets $A$ and $B$ is the set $A \times B$
of ordered pairs $(a, b)$ with $a \in A$ and $b \in B$. The product
preserves the source identity of each component: given a pair, one
can extract the first component (it lies in $A$) and the second
component (it lies in $B$). The two components do not merge into a
single value.

A measurement state space combining temperature and pressure is the
Cartesian product $T \times P$, where $T$ is the set of admissible
temperatures and $P$ is the set of admissible pressures. A specific
state is a pair $(t, p)$: a temperature value paired with a pressure
value. The state preserves both components; it does not combine them
into one undifferentiated number.

Given a state $(t, p)$, the temperature is recovered by the
projection $\pi_T : T \times P \to T$, $\pi_T(t, p) = t$. The
pressure is recovered by $\pi_P : T \times P \to P$, $\pi_P(t, p) =
p$. The projections are the algebraic mechanism for recovering the
components from the product. Without the projections, the product
would be a black box; with them, the components are accessible.

The chapter's second claim is that product structure is the
algebraic form of combination without collapse. The product
preserves both components; the projections recover them. Multiple
records combine via the product; the product's structure remembers
their identities.

The Cartesian product extends to arbitrarily many factors. The
product $A_1 \times A_2 \times \cdots \times A_n$ is the set of
tuples $(a_1, a_2, \ldots, a_n)$ with $a_i \in A_i$. Each
projection $\pi_i$ recovers the $i$-th component. The product
preserves all $n$ source identities.

Mathematically, the product is the categorical product in the
category of sets. The category-theoretic definition is: a product
of $A$ and $B$ is a set $P$ together with two maps $f_A : P \to A$
and $f_B : P \to B$, such that for any other set $C$ with maps
$g_A : C \to A$ and $g_B : C \to B$, there exists a unique map
$h : C \to P$ with $f_A \circ h = g_A$ and $f_B \circ h = g_B$.
This is the universal property of the product. The Cartesian
product $A \times B$ with its projections is the unique
(up to isomorphism) set satisfying this property.

In other categories, the product has different concrete forms. In
the category of groups, the product is the direct product, which
is a group with componentwise operation. In the category of
topological spaces, the product is the product topology. In the
category of measurable spaces, the product is the product
measurable space with the product $\sigma$-algebra. Each category
has its own product, but all share the universal property and the
component projections.

In Lean, the product type is \texttt{Prod $\alpha$ $\beta$}, which
is a structure with fields \texttt{fst} and \texttt{snd}. The
constructor takes two values and produces a pair; the projections
extract the components. The product type generalizes to dependent
sums (Sigma types) when the second component's type depends on
the first.

\begin{leanexcerpt}{Prod}
structure Prod ($\alpha$ : Type) ($\beta$ : Type) where
  fst : $\alpha$
  snd : $\beta$
\end{leanexcerpt}

The chapter's claim is that combination by product is the simplest
algebraic structure for multi-component records. Two records of
types $\alpha$ and $\beta$ combine into a record of type
\texttt{Prod $\alpha$ $\beta$}; the components are recoverable;
no information is lost.

The contrast is with combination by aggregation. Two integers
$5$ and $7$ can be combined into a sum $5 + 7 = 12$. The sum is
a single integer; the original components are not directly
recoverable (the sum could have been $4 + 8$, $3 + 9$, etc.).
The aggregation has lost the source identity.

Aggregation is appropriate for some purposes (computing totals,
averages, marginals) but not for others (preserving the audit
trail). The chapter's distinction is that the product structure
preserves identity; aggregation does not. Mathematical and
scientific records typically require identity preservation; the
product structure is the algebraic form for that requirement.

\section{Source / Encode / Execute}
\label{se:source-encode-execute}

A TeX source file, when compiled, produces a PDF. When the PDF is
opened, the viewer renders it on a screen. Three mathematical objects
are involved: the source (the TeX file as a string of characters); the
encoded form (the PDF as a sequence of bytes); the rendered output
(the visual image on the screen). The three objects are related but
distinct.

The source is what the author writes. The TeX file
\texttt{document.tex} is a text file: a sequence of characters
including TeX commands like \texttt{\textbackslash documentclass},
\texttt{\textbackslash section}, and the document's prose. The
characters are the mathematical content of the source.

The encoded form is the result of running the source through the
TeX compiler. The PDF \texttt{document.pdf} is a binary file: a
sequence of bytes encoding the typeset document. The bytes are
not the same as the characters of the source; they include
layout information, font instructions, page positions, and many
auxiliary structures.

The rendered output is what appears on the screen when the PDF is
opened. The viewer parses the PDF, computes the layout, draws the
glyphs at their positions, and displays the result. The visual
image is yet a third mathematical object: a function from screen
pixels to colors.

The three objects are not interchangeable. A reader cannot read
the source directly (it is full of TeX commands); a reader cannot
read the encoded PDF directly (the bytes are not human-readable);
the reader reads the rendered output. The rendering is the
admissible final form for human consumption. The source and the
encoded form are intermediate stages.

The chapter's third claim is that mathematical objects have stages
of realization. The source is the abstract specification; the
encoded form is the materialized instance; the executed output is
the rendered consequence. Each stage is a different mathematical
object; they are connected by transformations (compilation,
rendering); the chain preserves the document's identity through
the stages.

In Lean, this structure appears throughout. The Lean source code
\texttt{Example.lean} is the source. The elaborated terms (the
internal representation after type-checking) are the encoded
form. The compiled binary (the native code produced by Lean's
backend) is the executed form. Each stage is a different
mathematical object; the three are related by compilation and
elaboration.

\begin{leanexcerpt}{Encoding}
inductive Encoding
  | boot : Fact $\to$ Equivalation $\to$ Encoding
  | zero : Fact $\to$ Equivalation $\to$ Encoding $\to$ Encoding
  | one  : Fact $\to$ Equivalation $\to$ Equivalation
        $\to$ Encoding $\to$ Encoding $\to$ Encoding
\end{leanexcerpt}

The \texttt{Encoding} inductive in \texttt{Episode5.lean} carries
the staged representation directly. Each constructor corresponds
to a step in the digital cascade: \texttt{boot} for the initial
fact, \texttt{zero} and \texttt{one} for the bit positions that
extend a prior encoding into a longer one. The \texttt{Equivalation}
that each constructor carries is the bookkeeping that explains why
two stages can be treated as the same fact under transformation.

\texttt{SOURCE}, \texttt{EXECUTED}, and \texttt{Abstraction} are
the chapter's algebraic markers for the three stages. A record
in the source stage is a syntactic specification; a record in
the encoded stage is the elaborated form; a record in the
executed stage is the final output. The \texttt{Abstraction}
inductive certifies that a record has been correctly transformed
through the stages.

The chapter's claim is that the staged structure preserves
identity through transformations. The source's identity is
preserved in the encoded form (the same document content); the
encoded form's identity is preserved in the executed output
(the same visual layout). The three are distinct mathematical
objects, but they refer to the same document.

This is essential for the chapter's algebra. When two records are
combined, the combination should preserve the source identity of
each component. The staged structure ensures that the
combination's stages are tracked: if record A is in source stage
and record B is in encoded stage, the combination is in mixed
stage; the result must be tracked.

A practical example. A lambda calculus expression $(\lambda x. x +
1) 3$ is in source form. Beta reduction transforms it to $3 + 1$
(encoded). Arithmetic evaluation transforms it to $4$ (executed).
Each stage is distinct: $\lambda x. x + 1$ applied to $3$ is not
the same syntactic object as $3 + 1$, even though they evaluate
to the same value. The lambda calculus is a formal system that
distinguishes these stages explicitly.

The chapter's discipline is that combinations track the stages. A
record at the source stage combined with a record at the
executed stage cannot directly produce a result at either stage;
the combination is at the lowest common stage, which is source.
The result must be tracked through the subsequent stages to
reach the executed form.

\section{Unresolved States}
\label{se:unresolved-states}

A formal power series is an infinite sum $\sum_{n=0}^\infty a_n x^n$
where the coefficients $a_n$ are elements of a commutative ring (say,
the real numbers) and $x$ is a formal variable. The variable $x$ is
not assigned a specific value; it remains a placeholder. The series
is admissible as a formal mathematical object even though no value
has been substituted for $x$.

The series carries structure even unresolved. It can be added to
another formal power series (componentwise). It can be multiplied by
another series (Cauchy product). It can be differentiated (term by
term). It can be inverted (when the constant term is nonzero). The
operations are admissible at the level of the formal series, without
ever assigning $x$ a value.

The series has a convergence radius: the largest value $|x| = r$ for
which the series converges. For the series $\sum 1/n^2 \cdot x^n$, the
radius is $1$: the series converges for $|x| < 1$ and diverges for
$|x| > 1$. The radius is computable from the coefficients: $r =
1 / \limsup_{n \to \infty} |a_n|^{1/n}$ (Cauchy-Hadamard formula).

The formal series is the chapter's instance of an unresolved state.
The variable $x$ is an unresolved component; the series's structure
constrains what $x$ can be (within the convergence radius); the
operations on the series are admissible without resolving $x$.

The chapter's fourth claim is that unresolved states are admissible
mathematical objects. A record can carry an unresolved variable
indefinitely, as long as the structural constraints (the
coefficients, the convergence radius, the operations) are recorded.
The unresolved state is not an error; it is a deliberately
constrained placeholder.

\begin{leanexcerpt}{Jar}
inductive Jar
  | color         : Fact $\to$ Area $\to$ Jar
  | bang          : Fact $\to$ Jar $\to$ Jar
  | superposition : Fact $\to$ Jar $\to$ Jar $\to$ Jar
\end{leanexcerpt}

A \texttt{Jar} in \texttt{Episode4.lean} is an inductive with three
constructors: \texttt{color} (a base shape carrying a \texttt{Fact}
and an \texttt{Area}), \texttt{bang} (a single-step disturbance of
a prior \texttt{Jar}), and \texttt{superposition} (a layered pairing
of two prior \texttt{Jar}s under a single \texttt{Fact}). The
constructors do not assert a resolved value. They record the
admissible shape of the unresolved interior. A \texttt{Jar} can be
passed around, ordered against other \texttt{Jar}s, and combined
with them, all while leaving the interior unresolved.

The \texttt{MeesaProcess} typeclass handles \texttt{Jar}s:

\begin{leanexcerpt}{MeesaProcess}
structure MeesaProcess
    (Value : Type)
    (Carrier : CarrierProcess Value)
    [DISTINGUISHABLE Value Carrier]
    [ADMISSIBLE Value Carrier]
    [COUNTABLE Value Carrier]
    [ENCODED Value Carrier]
    [RESIDUE Value Carrier]
    [BINARY Value Carrier]
    [REPEATABLE Value Carrier]
    [NUMERIC Value Carrier]
    [REPRESENTABLE Value Carrier]
    [PHYSICAL Value Carrier]
    [COMPARABLE Value Carrier]
    [OBSERVED Value Carrier]
    [PRESENT Value Carrier]
    [MEASURABLE Value Carrier]
  where
  gauge_process : GaugeProcess Value Carrier
  concept       : Jar
  life_debt?    : Jar $\to$ Jar := ...
\end{leanexcerpt}

A \texttt{MeesaProcess} carries the project's full class cascade
through \texttt{MEASURABLE}, plus a \texttt{gauge\_process}, a
\texttt{concept} (the current \texttt{Jar}), and a function
\texttt{life\_debt?} that decides how a \texttt{Jar} is rewritten
under the process. The function is the algorithmic decision rule;
the cascade is the proof that the rule sits on every prior class
admissibly. The \texttt{GUNGAN} certificate confirms that the
process handles unresolved states admissibly: every unresolved
\texttt{Jar} is
either preserved or resolved according to the rule; no
\texttt{Jar} is silently mishandled.

A practical example. A clinical trial enrolls patients, administers
the experimental drug, and collects outcome data. The trial is
ongoing; the data are still being gathered. The primary endpoint
(say, two-year survival rate) is an unresolved state: the actual
outcomes have not all occurred yet. The trial's analysis must
treat the endpoint as unresolved; pretending the endpoint has a
specific value before the trial completes would be inadmissible.

When the trial completes, the endpoint is resolved: a specific
value is computed from the collected data. The unresolved state
transitions to resolved. The transition is admissible because
the data are now in. Before the transition, the analysis must
work with the unresolved state.

Mathematical record-keeping has the same discipline. An unresolved
state is admissible; it is not a defect or an error. The state is
constrained; the constraint allows admissible reasoning about
what the resolved state could be. Premature resolution
(assigning a value before the data warrant it) is the
inadmissible operation; honest preservation of the unresolved
state is the admissible operation.

The chapter's claim is that mathematical records routinely carry
unresolved states. Formal power series, free variables in
algebraic expressions, parameters in differential equations,
metavariables in type inference --- all are admissible mathematical
objects with unresolved components. The records' algebra includes
operations on the unresolved states without requiring their
resolution.

\section{Staging}
\label{se:staging}

The lambda calculus expression $(\lambda x. x + 1) \, 3$ is a program. The
expression is in its abstract form: a function ($\lambda x. x + 1$)
applied to an argument ($3$). The expression has not yet been
evaluated. The application is a syntactic construct; the evaluation
is a separate operation.

To evaluate the expression, apply the beta-reduction rule:
substitute the argument for the bound variable. $(\lambda x. x + 1) \,
3$ becomes $3 + 1$. The expression has been transformed by one beta
reduction step. The result $3 + 1$ is still a syntactic expression,
not a final value.

To complete the evaluation, apply the arithmetic rule: compute $3 + 1
= 4$. The expression has been transformed by one arithmetic step. The
result $4$ is the final value: an integer, not requiring further
evaluation.

Three stages are present. The original $(\lambda x. x + 1) \, 3$ is the
\emph{program}. The intermediate $3 + 1$ is the \emph{evaluated
application}. The final $4$ is the \emph{value}. The stages are
distinct: the program is a function and an argument; the
intermediate is a sum waiting to be computed; the value is a
specific integer. The transformations between stages are beta
reduction and arithmetic evaluation.

This is the chapter's fifth claim. Records have stages that
correspond to the algebraic operations performed on them. Each
stage is a different mathematical object; the transformations
between stages are well-defined operations. The records' algebra
includes both the operations and the staging.

Lambda calculus formalizes this. The expression $(\lambda x. x + 1)$ is
a function: a closed lambda term. The expression $3$ is a literal.
The expression $(\lambda x. x + 1) \, 3$ is an application: a function
applied to an argument. The expression $3 + 1$ is a different
application: addition applied to two literals. The expression $4$ is
a literal. Each is a distinct syntactic object.

Beta reduction is the rule: $(\lambda x. M) \, N$ reduces to $M[N/x]$
(the body $M$ with $N$ substituted for $x$). The rule is
deterministic when no two redexes interact; non-determinism appears
in cases like $(\lambda x. (\lambda y. M)) \, (\lambda z. N)$ where the
order of reductions matters. The Church-Rosser theorem says that, for
the simply typed lambda calculus, all reduction strategies converge to
the same normal form. The normal form is the unique value.

In Lean, the lambda calculus is the underlying calculus of the
type theory. Every term reduces to a normal form (when the type
theory is normalizing). The compiler's elaboration performs this
reduction during type checking. The reductions are the
transformations between stages; the normal form is the final
value.

\begin{leanexcerpt}{Equivalation}
inductive Equivalation
  | physics   : Fact $\to$ Jar $\to$ Equivalation
  | zero_like : Fact $\to$ Equivalation $\to$ Equivalation
  | one_like  : Prop $\to$ Equivalation $\to$ Equivalation
              $\to$ Equivalation
\end{leanexcerpt}

The \texttt{Equivalation} inductive in \texttt{Episode5.lean}
sits underneath \texttt{Encoding} and carries the explanation for
why two facts can be staged as the same value. Each constructor
records a different shape of equivalence: \texttt{physics} ties a
\texttt{Fact} to a \texttt{Jar}; \texttt{zero\_like} extends a
prior equivalation by a null step; \texttt{one\_like} pairs two
equivalations under a propositional witness. The
\texttt{CompiledProcess} structure certifies that the process has
progressed through the stages admissibly.

The chapter's claim is that staging is the algebraic discipline
of progressive realization. A record at the source stage is an
abstract specification. The encoded stage is the materialized
intermediate. The executed stage is the final value. The
algebra of records tracks the stage; operations admissible at
one stage may be inadmissible at another.

A practical example. A SQL query is a source: a logical
specification of what data to retrieve. The query plan is the
encoded form: an optimized procedure that the database engine
will execute. The result set is the executed value: the
specific rows returned. The three stages are tracked separately;
the database can reason about each stage's properties
(complexity, performance, expected size).

If the database treated all three stages as one, the
optimization opportunities would vanish (no way to reason about
the query plan separately from the result), and the auditing
trail would be lost (no way to reproduce the result if the data
changes). The staging is what enables the optimization and the
audit. The discipline is to track the stages explicitly.

The chapter closes here. Compatible records combine via products
without collapse. The product structure preserves source
identity. Source/encode/execute is the staging structure: each
record has stages of realization, and the transformations
between stages are admissible operations. Unresolved states
admit operations without requiring resolution. Staging tracks
the records' progressive realization.

\begin{bridge}

The chapter's question was how multiple records combine without
collapsing into one undifferentiated record.

The answer is the algebra of records. Compatibility constraints
determine which records can be combined; the product structure
preserves source identity; the source/encode/execute staging
tracks the records' progressive realization; unresolved states
admit operations without requiring resolution; staging tracks
the algebraic transformations between stages.

The chapter has not yet introduced distance. Two records can be
combined; the product structure preserves their identities; the
staging tracks their realization. But two records do not yet have
a notion of how close or how far they are from each other.
Chapter 5 takes up the question: how does a finite count of
distinctions become an admissible notion of distance?

\end{bridge}

\begin{coda}{The Stage Production Before Opening Night}

The theater has been preparing for three months. Opening night is
in two weeks. A contemporary staging of a classical Greek tragedy
is in technical rehearsal. Multiple production teams work in
parallel on different aspects of the same performance.

The records of the production are five distinct mathematical objects.

First, the script. The script is the source: the playwright's words,
the stage directions, the structure of acts and scenes. The script
has been worked on by the director and the dramaturg; it has been
edited, trimmed in places, expanded in others. The current version
is the rehearsal script: not the playwright's original, but the
production's specific text. The script is bound in three-ring
binders; each cast member has a copy with their lines marked.

Second, the blocking diagram. The blocking is the spatial
arrangement of the actors on the stage: where each character
stands at each moment, how they move during the scene, how they
exit. The blocking is recorded in the stage manager's prompt
book: a copy of the script with diagrams in the margins showing
the positions. The blocking is the script's encoded form for
performance: the script's words placed in physical space.

Third, the lighting cue sheet. The lighting designer has
prepared a list of cues: at each moment in the production, the
designer specifies which lamps are on, at what intensity, in what
color. The cue sheet is its own record, separate from the
script and the blocking, though it references both (each cue is
tied to a specific moment in the script). The cue sheet is what
the lighting board operator will run during the performance.

Fourth, the sound cue sheet. The sound designer has a parallel
list: at each moment, which sounds are played, at what volume,
in what panning. The sound cue sheet is structurally similar to
the lighting cue sheet but for the auditory dimension. It is also
a separate record.

Fifth, the rehearsal log. The stage manager keeps a log of each
rehearsal: what was rehearsed, who was present, what notes were
given, what issues were identified. The log is the production's
history: a record of how the production has evolved from the
script to its current state.

Each of the five records is mathematically distinct. The script is
a sequence of words. The blocking is a function from time to spatial
positions. The lighting cue sheet is a function from time to
lighting state. The sound cue sheet is a function from time to
sound state. The rehearsal log is a sequence of events in time.

The five records are also related. The script is the source; the
blocking is the script encoded in space; the lighting and sound
cue sheets are the script encoded in light and sound; the
rehearsal log is the history of these encodings being refined.
The relations among the records form the chapter's product
structure: the production is the product of the five records,
preserving each one's identity while assembling them into the
performance.

The product structure is what makes the production admissible.
Without it --- if the records were merged into one
undifferentiated production-blob --- the team would lose the
ability to revise specific aspects. The lighting designer could
not adjust the cue sheet without unmoving the script; the
costumer could not change a costume without affecting the
blocking. The records' separation is what enables the
production to be modified.

Compatibility constraints among the records are real. The
lighting cue sheet must reference moments that exist in the
script; a cue for ``the moment when the chorus enters'' is
meaningful only if the script has such a moment. The blocking
must place actors in positions where the stage permits movement;
a blocking that requires an actor to walk through a wall is
inadmissible. The compatibility is checked at the rehearsal
stage: technical rehearsals are when the records are tested
against each other, and incompatibilities surface for
resolution.

The production progresses through stages. The earliest stage is
the read-through: the cast reads the script aloud. The script is
in source stage; the cast's interpretation is in
\emph{abstraction} (in the sense of the chapter's typeclass):
the script's meaning is being assembled in the cast's
understanding. The next stage is blocking rehearsal: the
director places the actors in space. The script is being
encoded into the spatial form. The next is technical rehearsal:
the lighting, sound, and effects are integrated. The records
are being encoded into a unified performance plan. The next is
dress rehearsal: the costumes, makeup, and props are added.
The next is opening night: the executed performance, with the
audience present.

Each stage is a different mathematical object. The script alone
is one stage; the script plus blocking is another; the script
plus blocking plus tech is another; the full integrated
performance is the executed stage. The stages are distinct;
operations admissible at one stage may be inadmissible at another
(a major script revision is admissible during early rehearsal;
it is inadmissible the day before opening night).

The stage manager's logging discipline tracks all of this. Every
rehearsal note is recorded with the date, the stage of the
production, the participants, the issue, and the resolution. The
log is admissible only because of the chapter's algebra: each
note is a record; the notes combine via the rehearsal log's
product structure; the log preserves each note's identity; the
log tracks the production's progressive stages.

Unresolved states are common. A specific lighting cue is
``unresolved --- waiting for designer's decision.'' A specific
prop is ``unresolved --- waiting for substitute to arrive.'' A
specific blocking moment is ``unresolved --- to be revisited
after the actor returns from vacation.'' Each unresolved state
is admissible: the production proceeds with the unresolved
state carried explicitly. The stage manager's log marks each as
\texttt{UNRESOLVED} with a note about what is needed to resolve
it.

By opening night, almost all of the unresolved states must have
been resolved. The exception is items that are deliberately left
flexible (the actors' choices in moments of improvisation, the
audience's reaction times). These are admissible flexibilities,
not unresolved states; the difference is that the flexibilities
are part of the design, while unresolved states are pending
decisions.

When the curtain rises on opening night, the production is in
its executed stage. The script, the blocking, the lighting, the
sound, and the rehearsal log are all behind it; the audience
sees the result. The result is the product of the five records,
admissible because the records were compatible at every prior
stage. Without the product structure, the performance would be
a collection of independent elements; with it, the performance
is a unified artwork.

The five records do not collapse into the performance; they are
preserved as distinct documents. The script remains the
playwright's text plus the production's edits. The blocking
remains in the stage manager's prompt book. The cue sheets
remain in the designers' files. The rehearsal log remains in
the archive. The performance is what they together produce, but
each record is preserved separately for future productions or
revivals.

The chapter's structure appears throughout. The records are
compatible (each constrains the others' admissible structure).
The product structure preserves their identities. The
source/encode/execute staging tracks their progressive
realization. Unresolved states are admissible throughout the
rehearsal process. Staging tracks the algebraic transformations
from script to performance. None of the records is dissolved into
the performance; the performance is the product of records that
remain individually identifiable.

\end{coda}
